\clearpage
% 三类典型题型导航图
% 用法：% 三类典型题型导航图
% 用法：% 三类典型题型导航图
% 用法：\input{assets/tikz/R032_example_three_cases.tex}

\begin{center}
	\begin{tikzpicture}[
		x=1cm,
		y=1cm,
		>=Stealth,
		line cap=round,
		line join=round,
		every node/.style={font=\small},
		axis/.style={-Stealth, black!60, line width=0.65pt},
		panel/.style={rounded corners=7pt, line width=0.85pt}
	]
		\path[use as bounding box] (0,-0.05) rectangle (12.6,3.95);
		\node[font=\bfseries\large, text=CExample] at (6.3,3.62)
			{三类题型，一眼识别概率区域};

		% Panel 1
		\fill[green!3, rounded corners=7pt] (0.15,0.15) rectangle (4.05,3.25);
		\draw[green!50!black, panel] (0.15,0.15) rectangle (4.05,3.25);
		\node[font=\bfseries, text=CExample] at (2.10,2.92) {对称区间};
		\node[font=\scriptsize, text=black!65] at (2.10,2.58) {例 1：$40<X<60$};
		\draw[axis] (0.55,1.62) -- (3.65,1.62);
		\draw[green!65!black, line width=4.5pt] (1.12,1.62) -- (3.08,1.62);
		\foreach \x/\lab in {1.12/{40},2.10/{50},3.08/{60}}{
			\draw[black!65] (\x,1.50) -- (\x,1.74);
			\node[below=2pt, font=\scriptsize] at (\x,1.52) {$\lab$};
		}
		\node[font=\scriptsize, text=CExample] at (2.10,0.72)
			{$\mu\pm\sigma\Rightarrow68.27\%$};
		\node[font=\scriptsize, text=black!60] at (2.10,0.38)
			{直接用中央覆盖率};

		% Panel 2
		\fill[orange!3, rounded corners=7pt] (4.35,0.15) rectangle (8.25,3.25);
		\draw[orange!70!black, panel] (4.35,0.15) rectangle (8.25,3.25);
		\node[font=\bfseries, text=orange!80!black] at (6.30,2.92) {单侧尾部};
		\node[font=\scriptsize, text=black!65] at (6.30,2.58) {例 2：$X>70$};
		\draw[axis] (4.75,1.62) -- (7.85,1.62);
		\draw[orange!80!black, -Stealth, line width=4.5pt] (6.95,1.62) -- (7.76,1.62);
		\foreach \x/\lab in {5.35/{50},6.95/{70}}{
			\draw[black!65] (\x,1.50) -- (\x,1.74);
			\node[below=2pt, font=\scriptsize] at (\x,1.52) {$\lab$};
		}
		\node[font=\scriptsize, text=orange!80!black] at (6.30,0.72)
			{$Z>2\Rightarrow1-\Phi(2)$};
		\node[font=\scriptsize, text=black!60] at (6.30,0.38)
			{取补集或利用对称性};

		% Panel 3
		\fill[blue!3, rounded corners=7pt] (8.55,0.15) rectangle (12.45,3.25);
		\draw[blue!60!black, panel] (8.55,0.15) rectangle (12.45,3.25);
		\node[font=\bfseries, text=blue!70!black] at (10.50,2.92) {非对称区间};
		\node[font=\scriptsize, text=black!65] at (10.50,2.58) {例 3：$-1<Z<2$};
		\draw[axis] (8.95,1.62) -- (12.05,1.62);
		\draw[blue!65!black, line width=4.5pt] (9.48,1.62) -- (11.56,1.62);
		\foreach \x/\lab in {9.48/{-1},10.18/{0},11.56/{2}}{
			\draw[black!65] (\x,1.50) -- (\x,1.74);
			\node[below=2pt, font=\scriptsize] at (\x,1.52) {$\lab$};
		}
		\node[font=\scriptsize, text=blue!70!black] at (10.50,0.72)
			{$\Phi(2)-\Phi(-1)$};
		\node[font=\scriptsize, text=black!60] at (10.50,0.38)
			{分别标准化两个端点};
	\end{tikzpicture}
\end{center}


\begin{center}
	\begin{tikzpicture}[
		x=1cm,
		y=1cm,
		>=Stealth,
		line cap=round,
		line join=round,
		every node/.style={font=\small},
		axis/.style={-Stealth, black!60, line width=0.65pt},
		panel/.style={rounded corners=7pt, line width=0.85pt}
	]
		\path[use as bounding box] (0,-0.05) rectangle (12.6,3.95);
		\node[font=\bfseries\large, text=CExample] at (6.3,3.62)
			{三类题型，一眼识别概率区域};

		% Panel 1
		\fill[green!3, rounded corners=7pt] (0.15,0.15) rectangle (4.05,3.25);
		\draw[green!50!black, panel] (0.15,0.15) rectangle (4.05,3.25);
		\node[font=\bfseries, text=CExample] at (2.10,2.92) {对称区间};
		\node[font=\scriptsize, text=black!65] at (2.10,2.58) {例 1：$40<X<60$};
		\draw[axis] (0.55,1.62) -- (3.65,1.62);
		\draw[green!65!black, line width=4.5pt] (1.12,1.62) -- (3.08,1.62);
		\foreach \x/\lab in {1.12/{40},2.10/{50},3.08/{60}}{
			\draw[black!65] (\x,1.50) -- (\x,1.74);
			\node[below=2pt, font=\scriptsize] at (\x,1.52) {$\lab$};
		}
		\node[font=\scriptsize, text=CExample] at (2.10,0.72)
			{$\mu\pm\sigma\Rightarrow68.27\%$};
		\node[font=\scriptsize, text=black!60] at (2.10,0.38)
			{直接用中央覆盖率};

		% Panel 2
		\fill[orange!3, rounded corners=7pt] (4.35,0.15) rectangle (8.25,3.25);
		\draw[orange!70!black, panel] (4.35,0.15) rectangle (8.25,3.25);
		\node[font=\bfseries, text=orange!80!black] at (6.30,2.92) {单侧尾部};
		\node[font=\scriptsize, text=black!65] at (6.30,2.58) {例 2：$X>70$};
		\draw[axis] (4.75,1.62) -- (7.85,1.62);
		\draw[orange!80!black, -Stealth, line width=4.5pt] (6.95,1.62) -- (7.76,1.62);
		\foreach \x/\lab in {5.35/{50},6.95/{70}}{
			\draw[black!65] (\x,1.50) -- (\x,1.74);
			\node[below=2pt, font=\scriptsize] at (\x,1.52) {$\lab$};
		}
		\node[font=\scriptsize, text=orange!80!black] at (6.30,0.72)
			{$Z>2\Rightarrow1-\Phi(2)$};
		\node[font=\scriptsize, text=black!60] at (6.30,0.38)
			{取补集或利用对称性};

		% Panel 3
		\fill[blue!3, rounded corners=7pt] (8.55,0.15) rectangle (12.45,3.25);
		\draw[blue!60!black, panel] (8.55,0.15) rectangle (12.45,3.25);
		\node[font=\bfseries, text=blue!70!black] at (10.50,2.92) {非对称区间};
		\node[font=\scriptsize, text=black!65] at (10.50,2.58) {例 3：$-1<Z<2$};
		\draw[axis] (8.95,1.62) -- (12.05,1.62);
		\draw[blue!65!black, line width=4.5pt] (9.48,1.62) -- (11.56,1.62);
		\foreach \x/\lab in {9.48/{-1},10.18/{0},11.56/{2}}{
			\draw[black!65] (\x,1.50) -- (\x,1.74);
			\node[below=2pt, font=\scriptsize] at (\x,1.52) {$\lab$};
		}
		\node[font=\scriptsize, text=blue!70!black] at (10.50,0.72)
			{$\Phi(2)-\Phi(-1)$};
		\node[font=\scriptsize, text=black!60] at (10.50,0.38)
			{分别标准化两个端点};
	\end{tikzpicture}
\end{center}


\begin{center}
	\begin{tikzpicture}[
		x=1cm,
		y=1cm,
		>=Stealth,
		line cap=round,
		line join=round,
		every node/.style={font=\small},
		axis/.style={-Stealth, black!60, line width=0.65pt},
		panel/.style={rounded corners=7pt, line width=0.85pt}
	]
		\path[use as bounding box] (0,-0.05) rectangle (12.6,3.95);
		\node[font=\bfseries\large, text=CExample] at (6.3,3.62)
			{三类题型，一眼识别概率区域};

		% Panel 1
		\fill[green!3, rounded corners=7pt] (0.15,0.15) rectangle (4.05,3.25);
		\draw[green!50!black, panel] (0.15,0.15) rectangle (4.05,3.25);
		\node[font=\bfseries, text=CExample] at (2.10,2.92) {对称区间};
		\node[font=\scriptsize, text=black!65] at (2.10,2.58) {例 1：$40<X<60$};
		\draw[axis] (0.55,1.62) -- (3.65,1.62);
		\draw[green!65!black, line width=4.5pt] (1.12,1.62) -- (3.08,1.62);
		\foreach \x/\lab in {1.12/{40},2.10/{50},3.08/{60}}{
			\draw[black!65] (\x,1.50) -- (\x,1.74);
			\node[below=2pt, font=\scriptsize] at (\x,1.52) {$\lab$};
		}
		\node[font=\scriptsize, text=CExample] at (2.10,0.72)
			{$\mu\pm\sigma\Rightarrow68.27\%$};
		\node[font=\scriptsize, text=black!60] at (2.10,0.38)
			{直接用中央覆盖率};

		% Panel 2
		\fill[orange!3, rounded corners=7pt] (4.35,0.15) rectangle (8.25,3.25);
		\draw[orange!70!black, panel] (4.35,0.15) rectangle (8.25,3.25);
		\node[font=\bfseries, text=orange!80!black] at (6.30,2.92) {单侧尾部};
		\node[font=\scriptsize, text=black!65] at (6.30,2.58) {例 2：$X>70$};
		\draw[axis] (4.75,1.62) -- (7.85,1.62);
		\draw[orange!80!black, -Stealth, line width=4.5pt] (6.95,1.62) -- (7.76,1.62);
		\foreach \x/\lab in {5.35/{50},6.95/{70}}{
			\draw[black!65] (\x,1.50) -- (\x,1.74);
			\node[below=2pt, font=\scriptsize] at (\x,1.52) {$\lab$};
		}
		\node[font=\scriptsize, text=orange!80!black] at (6.30,0.72)
			{$Z>2\Rightarrow1-\Phi(2)$};
		\node[font=\scriptsize, text=black!60] at (6.30,0.38)
			{取补集或利用对称性};

		% Panel 3
		\fill[blue!3, rounded corners=7pt] (8.55,0.15) rectangle (12.45,3.25);
		\draw[blue!60!black, panel] (8.55,0.15) rectangle (12.45,3.25);
		\node[font=\bfseries, text=blue!70!black] at (10.50,2.92) {非对称区间};
		\node[font=\scriptsize, text=black!65] at (10.50,2.58) {例 3：$-1<Z<2$};
		\draw[axis] (8.95,1.62) -- (12.05,1.62);
		\draw[blue!65!black, line width=4.5pt] (9.48,1.62) -- (11.56,1.62);
		\foreach \x/\lab in {9.48/{-1},10.18/{0},11.56/{2}}{
			\draw[black!65] (\x,1.50) -- (\x,1.74);
			\node[below=2pt, font=\scriptsize] at (\x,1.52) {$\lab$};
		}
		\node[font=\scriptsize, text=blue!70!black] at (10.50,0.72)
			{$\Phi(2)-\Phi(-1)$};
		\node[font=\scriptsize, text=black!60] at (10.50,0.38)
			{分别标准化两个端点};
	\end{tikzpicture}
\end{center}


\begin{examplebox}
	\textbf{\textcolor{CExample}{例 1（识别方差与标准差）}}

	\textbf{\textcolor{CExample}{题目：}}
	设 $X\sim N(50,100)$，求 $P(40<X<60)$。

	\textbf{\textcolor{CExample}{解题步骤：}}
	\begin{description}[style=nextline, leftmargin=2.8em, labelsep=0.5em, itemsep=0.45em]
		\item[\textbf{第一步（读取参数）：}]
			第二个参数是方差，故 $\mu=50$，$\sigma=\sqrt{100}=10$。
		\item[\textbf{第二步（识别区间）：}]
			$40=\mu-\sigma$，$60=\mu+\sigma$，所以所求事件为 $|X-\mu|<\sigma$。
		\item[\textbf{第三步（套用结论）：}]
			由中央覆盖率公式，
			\[
				P(40<X<60)=2\Phi(1)-1\approx0.6827.
			\]
	\end{description}
	\textbf{\textcolor{CExample}{关键结论：}}
	\textbf{$P(40<X<60)\approx0.6827$。}

\end{examplebox}

\clearpage
% 例 2：单侧尾部概率图
% 用法：% 例 2：单侧尾部概率图
% 用法：% 例 2：单侧尾部概率图
% 用法：\input{assets/tikz/R032_example_right_tail.tex}

\begin{center}
	\begin{tikzpicture}[
		x=1.52cm,
		y=4.6cm,
		>=Stealth,
		line cap=round,
		line join=round,
		every node/.style={font=\small},
		axis/.style={-Stealth, black!60, line width=0.7pt}
	]
		\path[use as bounding box] (-3.85,-0.24) rectangle (3.85,0.55);
		\fill[orange!35]
			(2,0) -- plot[domain=2:3.65, samples=55] (\x,{0.4*exp(-\x*\x/2)}) -- (3.65,0) -- cycle;
		\draw[axis] (-3.65,0) -- (3.66,0) node[right] {$z$};
		\draw[orange!80!black, line width=1.05pt]
			plot[domain=-3.55:3.55, samples=150] (\x,{0.4*exp(-\x*\x/2)});
		\draw[dashed, orange!85!black, line width=0.8pt] (2,0)--(2,0.31);
		\foreach \x/\lab in {0/{0},2/{2}}{
			\draw[black!55] (\x,-0.012)--(\x,0.012);
			\node[below=1pt] at (\x,-0.004) {$\lab$};
		}
		\node[font=\bfseries\large, text=orange!80!black] at (0,0.49)
			{单侧尾部：先定位临界点，再取补集};
		\node[
			draw=orange!75!black,
			fill=orange!4,
			rounded corners=5pt,
			inner xsep=8pt,
			inner ysep=3pt,
			text=orange!80!black
		] at (2.75,0.19) {$1-\Phi(2)\approx0.0228$};
		\node[
			draw=black!35,
			fill=white,
			rounded corners=5pt,
			inner xsep=7pt,
			inner ysep=3pt,
			font=\scriptsize,
			text=black!65
		] at (-2.15,0.18) {左侧累计面积 $\Phi(2)\approx0.9772$};
		\node[font=\scriptsize, text=black!65] at (0,-0.18)
			{$70=\mu+2\sigma\ \Longrightarrow\ P(X>70)=P(Z>2)$};
	\end{tikzpicture}
\end{center}


\begin{center}
	\begin{tikzpicture}[
		x=1.52cm,
		y=4.6cm,
		>=Stealth,
		line cap=round,
		line join=round,
		every node/.style={font=\small},
		axis/.style={-Stealth, black!60, line width=0.7pt}
	]
		\path[use as bounding box] (-3.85,-0.24) rectangle (3.85,0.55);
		\fill[orange!35]
			(2,0) -- plot[domain=2:3.65, samples=55] (\x,{0.4*exp(-\x*\x/2)}) -- (3.65,0) -- cycle;
		\draw[axis] (-3.65,0) -- (3.66,0) node[right] {$z$};
		\draw[orange!80!black, line width=1.05pt]
			plot[domain=-3.55:3.55, samples=150] (\x,{0.4*exp(-\x*\x/2)});
		\draw[dashed, orange!85!black, line width=0.8pt] (2,0)--(2,0.31);
		\foreach \x/\lab in {0/{0},2/{2}}{
			\draw[black!55] (\x,-0.012)--(\x,0.012);
			\node[below=1pt] at (\x,-0.004) {$\lab$};
		}
		\node[font=\bfseries\large, text=orange!80!black] at (0,0.49)
			{单侧尾部：先定位临界点，再取补集};
		\node[
			draw=orange!75!black,
			fill=orange!4,
			rounded corners=5pt,
			inner xsep=8pt,
			inner ysep=3pt,
			text=orange!80!black
		] at (2.75,0.19) {$1-\Phi(2)\approx0.0228$};
		\node[
			draw=black!35,
			fill=white,
			rounded corners=5pt,
			inner xsep=7pt,
			inner ysep=3pt,
			font=\scriptsize,
			text=black!65
		] at (-2.15,0.18) {左侧累计面积 $\Phi(2)\approx0.9772$};
		\node[font=\scriptsize, text=black!65] at (0,-0.18)
			{$70=\mu+2\sigma\ \Longrightarrow\ P(X>70)=P(Z>2)$};
	\end{tikzpicture}
\end{center}


\begin{center}
	\begin{tikzpicture}[
		x=1.52cm,
		y=4.6cm,
		>=Stealth,
		line cap=round,
		line join=round,
		every node/.style={font=\small},
		axis/.style={-Stealth, black!60, line width=0.7pt}
	]
		\path[use as bounding box] (-3.85,-0.24) rectangle (3.85,0.55);
		\fill[orange!35]
			(2,0) -- plot[domain=2:3.65, samples=55] (\x,{0.4*exp(-\x*\x/2)}) -- (3.65,0) -- cycle;
		\draw[axis] (-3.65,0) -- (3.66,0) node[right] {$z$};
		\draw[orange!80!black, line width=1.05pt]
			plot[domain=-3.55:3.55, samples=150] (\x,{0.4*exp(-\x*\x/2)});
		\draw[dashed, orange!85!black, line width=0.8pt] (2,0)--(2,0.31);
		\foreach \x/\lab in {0/{0},2/{2}}{
			\draw[black!55] (\x,-0.012)--(\x,0.012);
			\node[below=1pt] at (\x,-0.004) {$\lab$};
		}
		\node[font=\bfseries\large, text=orange!80!black] at (0,0.49)
			{单侧尾部：先定位临界点，再取补集};
		\node[
			draw=orange!75!black,
			fill=orange!4,
			rounded corners=5pt,
			inner xsep=8pt,
			inner ysep=3pt,
			text=orange!80!black
		] at (2.75,0.19) {$1-\Phi(2)\approx0.0228$};
		\node[
			draw=black!35,
			fill=white,
			rounded corners=5pt,
			inner xsep=7pt,
			inner ysep=3pt,
			font=\scriptsize,
			text=black!65
		] at (-2.15,0.18) {左侧累计面积 $\Phi(2)\approx0.9772$};
		\node[font=\scriptsize, text=black!65] at (0,-0.18)
			{$70=\mu+2\sigma\ \Longrightarrow\ P(X>70)=P(Z>2)$};
	\end{tikzpicture}
\end{center}


\begin{examplebox}
	\textbf{\textcolor{CExample}{例 2（单侧尾部）}}

	\textbf{\textcolor{CExample}{题目：}}
	设 $X\sim N(50,100)$，求 $P(X>70)$。

	\textbf{\textcolor{CExample}{解题步骤：}}
	\begin{description}[style=nextline, leftmargin=2.8em, labelsep=0.5em, itemsep=0.45em]
		\item[\textbf{第一步（识别临界点）：}]
			由 $\mu=50$，$\sigma=10$，有 $70=\mu+2\sigma$。
		\item[\textbf{第二步（求单侧尾部）：}]
			由补集公式，
			\[
				\begin{aligned}
					P(X>70)
					&=P(Z>2)=1-\Phi(2)\\
					&=0.0227501\ldots\approx0.0228.
				\end{aligned}
			\]
			也可用 $2\sigma$ 覆盖概率近似计算为 $(1-0.9545)/2=0.02275$。
	\end{description}
	\textbf{\textcolor{CExample}{关键结论：}}
	\textbf{$P(X>70)\approx0.0228$。}

\end{examplebox}

\clearpage
% 例 3：非对称区间概率图
% 用法：% 例 3：非对称区间概率图
% 用法：% 例 3：非对称区间概率图
% 用法：\input{assets/tikz/R032_example_asymmetric_interval.tex}

\begin{center}
	\begin{tikzpicture}[
		x=1.52cm,
		y=4.6cm,
		>=Stealth,
		line cap=round,
		line join=round,
		every node/.style={font=\small},
		axis/.style={-Stealth, black!60, line width=0.7pt}
	]
		\path[use as bounding box] (-3.85,-0.24) rectangle (3.85,0.55);
		\fill[blue!24]
			(-1,0) -- plot[domain=-1:2, samples=90] (\x,{0.4*exp(-\x*\x/2)}) -- (2,0) -- cycle;
		\draw[axis] (-3.65,0) -- (3.66,0) node[right] {$z$};
		\draw[blue!65!black, line width=1.05pt]
			plot[domain=-3.55:3.55, samples=150] (\x,{0.4*exp(-\x*\x/2)});
		\foreach \x/\lab in {-1/{-1},0/{0},2/{2}}{
			\draw[black!55] (\x,-0.012)--(\x,0.012);
			\node[below=1pt] at (\x,-0.004) {$\lab$};
		}
		\draw[dashed, blue!70!black, line width=0.8pt] (-1,0)--(-1,0.26);
		\draw[dashed, blue!70!black, line width=0.8pt] (2,0)--(2,0.09);
		\node[font=\bfseries\large, text=blue!70!black] at (0,0.49)
			{非对称区间：两个端点分别标准化};
		\node[
			draw=blue!65!black,
			fill=white,
			rounded corners=5pt,
			inner xsep=8pt,
			inner ysep=3pt,
			text=blue!70!black
		] at (0,0.23) {$\Phi(2)-\Phi(-1)\approx0.8186$};
		\node[font=\scriptsize, text=black!65] at (0,-0.18)
			{$(-1,8)\xrightarrow{\ z=(x-2)/3\ }(-1,2)$};
	\end{tikzpicture}
\end{center}


\begin{center}
	\begin{tikzpicture}[
		x=1.52cm,
		y=4.6cm,
		>=Stealth,
		line cap=round,
		line join=round,
		every node/.style={font=\small},
		axis/.style={-Stealth, black!60, line width=0.7pt}
	]
		\path[use as bounding box] (-3.85,-0.24) rectangle (3.85,0.55);
		\fill[blue!24]
			(-1,0) -- plot[domain=-1:2, samples=90] (\x,{0.4*exp(-\x*\x/2)}) -- (2,0) -- cycle;
		\draw[axis] (-3.65,0) -- (3.66,0) node[right] {$z$};
		\draw[blue!65!black, line width=1.05pt]
			plot[domain=-3.55:3.55, samples=150] (\x,{0.4*exp(-\x*\x/2)});
		\foreach \x/\lab in {-1/{-1},0/{0},2/{2}}{
			\draw[black!55] (\x,-0.012)--(\x,0.012);
			\node[below=1pt] at (\x,-0.004) {$\lab$};
		}
		\draw[dashed, blue!70!black, line width=0.8pt] (-1,0)--(-1,0.26);
		\draw[dashed, blue!70!black, line width=0.8pt] (2,0)--(2,0.09);
		\node[font=\bfseries\large, text=blue!70!black] at (0,0.49)
			{非对称区间：两个端点分别标准化};
		\node[
			draw=blue!65!black,
			fill=white,
			rounded corners=5pt,
			inner xsep=8pt,
			inner ysep=3pt,
			text=blue!70!black
		] at (0,0.23) {$\Phi(2)-\Phi(-1)\approx0.8186$};
		\node[font=\scriptsize, text=black!65] at (0,-0.18)
			{$(-1,8)\xrightarrow{\ z=(x-2)/3\ }(-1,2)$};
	\end{tikzpicture}
\end{center}


\begin{center}
	\begin{tikzpicture}[
		x=1.52cm,
		y=4.6cm,
		>=Stealth,
		line cap=round,
		line join=round,
		every node/.style={font=\small},
		axis/.style={-Stealth, black!60, line width=0.7pt}
	]
		\path[use as bounding box] (-3.85,-0.24) rectangle (3.85,0.55);
		\fill[blue!24]
			(-1,0) -- plot[domain=-1:2, samples=90] (\x,{0.4*exp(-\x*\x/2)}) -- (2,0) -- cycle;
		\draw[axis] (-3.65,0) -- (3.66,0) node[right] {$z$};
		\draw[blue!65!black, line width=1.05pt]
			plot[domain=-3.55:3.55, samples=150] (\x,{0.4*exp(-\x*\x/2)});
		\foreach \x/\lab in {-1/{-1},0/{0},2/{2}}{
			\draw[black!55] (\x,-0.012)--(\x,0.012);
			\node[below=1pt] at (\x,-0.004) {$\lab$};
		}
		\draw[dashed, blue!70!black, line width=0.8pt] (-1,0)--(-1,0.26);
		\draw[dashed, blue!70!black, line width=0.8pt] (2,0)--(2,0.09);
		\node[font=\bfseries\large, text=blue!70!black] at (0,0.49)
			{非对称区间：两个端点分别标准化};
		\node[
			draw=blue!65!black,
			fill=white,
			rounded corners=5pt,
			inner xsep=8pt,
			inner ysep=3pt,
			text=blue!70!black
		] at (0,0.23) {$\Phi(2)-\Phi(-1)\approx0.8186$};
		\node[font=\scriptsize, text=black!65] at (0,-0.18)
			{$(-1,8)\xrightarrow{\ z=(x-2)/3\ }(-1,2)$};
	\end{tikzpicture}
\end{center}


\begin{examplebox}
	\textbf{\textcolor{CExample}{例 3（非对称区间）}}

	\textbf{\textcolor{CExample}{题目：}}
	设 $X\sim N(2,9)$，求 $P(-1<X<8)$。

	\textbf{\textcolor{CExample}{解题步骤：}}
	\begin{description}[style=nextline, leftmargin=2.8em, labelsep=0.5em, itemsep=0.45em]
		\item[\textbf{第一步（标准化两个端点）：}]
			$\mu=2$，$\sigma=3$，故
			\[
				\frac{-1-2}{3}=-1,
				\qquad
				\frac{8-2}{3}=2.
			\]
		\item[\textbf{第二步（用分布函数作差）：}]
			原区间不关于 $\mu$ 对称，不能直接套用某一个 $k\sigma$ 覆盖率。应计算
			\[
				\begin{aligned}
					P(-1<X<8)
					&=P(-1<Z<2)\\
					&=\Phi(2)-\Phi(-1)\\
					&=\Phi(2)+\Phi(1)-1\\
					&=0.8185946\ldots\approx0.8186.
				\end{aligned}
			\]
	\end{description}
	\textbf{\textcolor{CExample}{关键结论：}}
	\textbf{$P(-1<X<8)\approx0.8186$。}
\end{examplebox}
